\section{The gradient on the conditioned edge measure}
\label{sec:gradient}

The variance estimate of Proposition~\ref{prop:variance-main} says that a
function of bounded rescaled energy is almost determined by a sufficiently
fine cap observation. We now identify the derivatives of the resulting
functions of the edge coordinates. There are two issues. The measure
$\Pi$ is supported on a renormalized quadratic constraint, so its Sobolev
domain must first be specified. We then show that the slopes of the
observation smoothings converge to derivatives in this domain. The energy
bound obtained here has a finite constant; Section~\ref{sec:microscopic}
will determine the exact coefficient.

Throughout the first part of this section, fix a realization of the edge
parameters for which Proposition~\ref{prop:edge-conditioned-law} holds.
Thus $l_a=(g_a+b)^{-1}>0$, $\sum_a l_a=\infty$, and
$\sum_a l_a^2<\infty$. All assertions about $\Pi$ below are deterministic
for this realization.

\subsection{Angular derivatives and the form domain}

A \emph{smooth cylinder function} is a smooth function of finitely many
coordinates. We initially take it compactly supported in those coordinates.
For a skew-symmetric matrix $A$ with finitely many nonzero entries, set
$\mathcal R_A=(Ax)\cdot\partial_x$. We say that
$(F,G)\in L^2(\Pi)\times L^2(\Pi;\ell^2)$ satisfies the angular weak
identities if
\begin{equation}
 \int \psi\,(Ax)\cdot G\,d\Pi
 =-\int F\left\{\mathcal R_A\psi
       -\psi\sum_a (Ax)_a(g_a+b)x_a\right\}\,d\Pi
 \label{eq:source-23}
\end{equation}
for every such $A$ and every smooth compact cylinder test $\psi$ whose
coordinate set contains the indices of $A$. Finite rotations preserve the
quadratic constraint. Their action on the product Gaussian density gives
the second term in braces; the contribution involving $b$ cancels by
skew-symmetry. Consequently, smooth cylinders satisfy
\eqref{eq:source-23} with their Euclidean gradients.

\begin{proposition}[The conditioned gradient]
\label{prop:gradient-domain}
The Euclidean gradient on smooth compact cylinder functions is closable
from $L^2(\Pi)$ to $L^2(\Pi;\ell^2)$. Its closed graph consists exactly
of the pairs satisfying \eqref{eq:source-23}. In particular, the vector
$G$ in that identity is unique; we denote it by $\nabla F$.

The closed form
\[
 \mathcal E_\Pi(F)=\int|\nabla F|^2\,d\Pi
\]
has kernel equal to the constants. Each coordinate $x_a$ belongs to its
domain and satisfies $\nabla x_a=e_a$, so that
$\mathcal E_\Pi(x_a)=1$.
\end{proposition}

We first verify closability and then prove the density assertion, which is
the substantial part of the proposition. By the polar disintegration in
Proposition~\ref{prop:edge-conditioned-law}, the law of the coordinates
outside a sufficiently large fixed initial block is absolutely continuous with respect to
their product Gaussian law. Under that product law,
\[
 \sum_{a\le N}x_a^2
 =\sum_{a\le N}l_a+\sum_{a\le N}(x_a^2-l_a)\longrightarrow\infty
 \qquad\text{almost surely}.
\]
Thus $x\notin\ell^2$ also holds $\Pi$-almost surely. If $G$ and $G'$ both
satisfy \eqref{eq:source-23} for the same $F$, the coordinate rotations
give
\[
 x_i(G_j-G'_j)=x_j(G_i-G'_i)\qquad\text{for all }i,j,
\]
outside one null set. Their difference is therefore proportional to $x$
and must vanish because it belongs to $\ell^2$. The angular graph is
closed: every test in \eqref{eq:source-23} has bounded coefficients on its
finite coordinate support, so the identity passes to strong limits in
$L^2(\Pi)\times L^2(\Pi;\ell^2)$. A closed graph containing the cylinder
gradient and having a unique vector component proves closability.

\paragraph{Reduction to a product chart.}
Let $(F,G)$ satisfy \eqref{eq:source-23}. Conditional on the complement
of a finite block, its coordinates lie on a sphere of fixed radius with
a smooth positive tilted density. Disintegrating
\eqref{eq:source-23}, including tests in outside coordinates, gives the
ordinary tangential weak identities on almost every such sphere. The
rotation fields span its tangent spaces. The usual weak chain and
product rules therefore apply.

Smooth truncation first reduces to bounded $F$, with convergence of
both $F$ and $G$. Fix an initial block of size $m>3$, and write its
radius as $R$. We may next localize to an annulus
$\epsilon<R<M$. Near zero, the marginal density bound of
Proposition~\ref{prop:edge-conditioned-law} gives
$\Pi(R\le\epsilon)=O(\epsilon^m)$. A cutoff with derivative
$O(\epsilon^{-1})$ thus has energy $O(\epsilon^{m-2})$ when multiplied
by bounded $F$. At infinity, a cutoff with derivative $O(M^{-1})$
has vanishing cost. The corresponding products with $G$ converge by
dominated convergence. It suffices to approximate each localized pair.

Write the first block as $R\vartheta$, with
$\vartheta\in S^{m-1}$, and put $y=(x_i)_{i>m}$. The constraint gives
\[
 R^2=C-S_y,
 \qquad C=D(b)+\sum_{i\le m}l_i,
 \qquad S_y=\sum_{i>m}(y_i^2-l_i).
\]
Relative to uniform angle and the product Gaussian law of $y$, the
conditioned measure has density proportional to
\[
 \mathbf1_{\{S_y<C\}}R^{m-2}
 \exp\left(-\frac{R^2}{2}
                   \sum_{j\le m}\frac{\vartheta_j^2}{l_j}\right).
\]
On the chosen annulus this density is bounded above and below by
positive constants. Decompose the first block of $G$ into radial and
transverse components $G_R,G_\vartheta$, and let $f$ be the expression
of the localized $F$ in this chart. With $a=-G_R/(2R)$, the finite-block
identities give
\begin{equation}
 \nabla_\vartheta f=R G_\vartheta,
 \qquad
 \partial_{y_i}f-2y_i a=G_i\quad(i>m).
 \label{eq:source-24}
\end{equation}
To justify the weak tail identity, condition on the coordinates outside
the union of the first block and $\{i\}$. On this $(m+1)$-coordinate
sphere, varying $y_i$ at fixed angle has tangent vector
\[
 V_i=e_i-\frac{y_i}{R^2}\sum_{j\le m}x_j e_j.
\]
It is a smooth linear combination of the $i$--$j$ rotation fields on
$R>0$, and its pairing with $G$ is
$G_i-(y_i/R)G_R=G_i+2y_i a$. The conditional density is smooth and positive on
each compact chart with $R>0$, so multiplying tests by its reciprocal
converts the tangential weak identity to the ordinary weak coordinate
identity. Disintegration then gives \eqref{eq:source-24} in the product
chart, with Gaussian integration by parts for Gaussian tests. The radial
cutoff permits extension by zero without a boundary term. The quantities
$f$, $a$, $\nabla_\vartheta f$, and the compensated derivative vector
$(\partial_{y_i}f-2y_i a)_{i>m}$ are square integrable for angle times
product Gaussian measure.

The domain question has now become an approximation problem for the
compensated tail derivatives in \eqref{eq:source-24}. The following
lemma supplies exactly this approximation; the angular variable will
be handled afterward by ordinary spherical harmonics.

\begin{lemma}[Approximation with a compensated derivative]
\label{lem:gradient-compensated}
Let $\gamma=\bigotimes_{i\ge1}\mathsf N(0,l_i)$, where $l_i>0$,
$\sum_i l_i=\infty$ and $\sum_i l_i^2<\infty$, and let
$S_y=\sum_i(y_i^2-l_i)$. Suppose that $f,a\in L^2(\gamma)$ and
$G\in L^2(\gamma;\ell^2)$ satisfy, coordinatewise in the weak sense,
\[
 \partial_{y_i}f-2y_i a=G_i.
\]
Then $(f,a,G)$ is a limit in these three spaces of triples
\[
 \left(P,\partial_s P,
       \bigl(\partial_{y_i}P\bigr)_i\right)\Big|_{s=S_y},
\]
where $P=P(s,y_1,\ldots,y_N)$ is a polynomial. The derivatives in this
display treat $s$ and the individual $y_i$ as independent variables.
\end{lemma}

\begin{proof}
Expand $f$ and $a$ in the normalized product Hermite basis. We index the
basis by finitely supported multi-indices $\beta$ and write the
coefficients as $f_\beta,a_\beta$. Gaussian integration by parts gives
\begin{equation}
 (G_i)_\beta
 =\frac{\sqrt{\beta_i+1}}{\sqrt{l_i}}f_{\beta+e_i}
 -2\sqrt{l_i}\left(
      \sqrt{\beta_i}\,a_{\beta-e_i}
      +\sqrt{\beta_i+1}\,a_{\beta+e_i}\right).
 \label{eq:source-25}
\end{equation}
The formula is valid for weak derivatives: one first uses finite-coordinate
Hermite tests with cutoffs, then removes the cutoffs by Cauchy--Schwarz.

\paragraph{Reduction to one chaos grade.}
Group the unknowns as pairs $(f^{(d)},a^{(d-2)})$, where the superscript
denotes Hermite chaos order and negative orders are omitted. The derivative
of $f^{(d)}$ and the creation part of $2y_i a^{(d-2)}$ both have order
$d-1$ and must be kept together. The remaining, annihilation part of
$2y_i a^{(d)}$ has summed norm at most
\[
 \left(4\sum_{i,\beta}l_i(\beta_i+1)
                  |a_{\beta+e_i}|^2\right)^{1/2}
 \le C\sqrt d\,\|a^{(d)}\|_2.
\]
Here $\sup_i l_i<\infty$ follows from $\sum_i l_i^2<\infty$.
It follows from \eqref{eq:source-25} that each fixed-grade pair has a
square-summable compensated derivative.

These pairs also approximate the full triple in its graph norm. To see
this without separating divergent creation terms, multiply the grade-$d$
pair by $e^{-sd}$, with $s>0$. Relative to damping the output of
\eqref{eq:source-25} by $e^{-sd}$ in order $d-1$, the only discrepancy
is the annihilation term, with multiplier
$e^{-s(d+2)}-e^{-sd}$. Its squared norm is bounded by
\[
 C s^2\sum_d d e^{-2sd}\|a^{(d)}\|_2^2
 \le C s\,\|a\|_2^2,
\]
which tends to zero. For fixed $s$, truncating the damped grades then
converges in all three norms. It remains to approximate one grade.

\paragraph{The array represented by a fixed grade.}
Fix $d$ and the finite set $O$ of coordinates with odd multiplicity.
The relevant coefficients of $f$ have indices
$\mathbf1_O+2j$, where $j$ is a count vector with
$|j|=h=(d-|O|)/2$; those of $a$ use $|j|=h-1$.
Different parity sets are orthogonal in the norms of the common
derivative--creation terms, so finitely many such sets suffice for
approximation. If $h=0$, only a finite-coordinate monomial is involved.
Suppose henceforth that $h\ge1$.

Set
\[
 c(j)=\prod_i l_i^{j_i}\frac{\sqrt{(2j_i)!}}{j_i!},
 \qquad
 f_{\mathbf1_O+2j}=c(j)u(j),
 \qquad
 a_{\mathbf1_O+2j}=c(j)v(j).
\]
View $u$ and $v$ as symmetric arrays on $h$ and $h-1$ ordered places.
Give each place mass $\mu_i=l_i^2$. The coefficient norms are equivalent
to the resulting product $L^2$ norms: the ordered-place multiplicities
and the displayed factorials have ratios bounded above and below for
fixed $h$ and $O$. This comparison includes repeated indices.

For $i\notin O$, the identity
\[
 \frac{\sqrt{2j_i+2}}{\sqrt{l_i}}c(j+e_i)
 =2\sqrt{l_i}\sqrt{2j_i+1}\,c(j)
\]
shows that the common term in \eqref{eq:source-25} is exactly a multiple
of $u(j+e_i)-v(j)$. Its norm is therefore equivalent, again with
fixed-grade constants, to the sum over places of
\[
 \int \sum_i l_i\,
       |u(i,i_2,\ldots,i_h)-v(i_2,\ldots,i_h)|^2
       \,d\mu(i_2)\cdots d\mu(i_h).
\]
The finitely many indices in $O$ contribute bounded operators in the
coefficient norms, as does the annihilation term at this fixed grade.
We have thus reduced the approximation problem to symmetric arrays
with a common one-place trace $v$.

\paragraph{The sum of the one-place forms.}
For one place, take the base Hilbert space $\mathfrak H=\ell^2(\mu)$.
We include the scalar trace in the form norm by defining
\[
 \begin{split}
 D(q)&=\{u\in\mathfrak H:\text{there exists }v\in\mathbb R
                       \text{ with }u-v\in\ell^2(l_i)\},\\
 q(u)&=|v|^2+\sum_i l_i|u_i-v|^2.
 \end{split}
\]
The trace $v$ is unique because $\sum_i l_i=\infty$. This is a densely
defined closed form. Indeed, finite-support sequences belong to $D(q)$;
and a Cauchy sequence in the form norm has convergent traces and
convergent deviations, whose coordinatewise limits identify the trace
and deviation of its limit in $\mathfrak H$. Replacing the tail of $u$
by its trace proves that eventually constant sequences are a form core.

On $\mathfrak H^{\otimes h}$ take the \emph{sum} of the $h$ coordinate
forms. By the slice description of each coordinate form, its norm on
symmetric arrays is precisely, up to fixed factors, the norm of $u$,
the norm of $v$, and the integrated deviations just obtained. In
particular, it requires $v$ to belong to the base $(h-1)$-place Hilbert
space; it requires no further form regularity of $v$.

For completeness, let $A_q$ be the nonnegative selfadjoint operator
associated with $q$. The commuting projections
$\mathbf1_{[0,M]}(A_q)^{\otimes h}$ converge to the identity in the
norm of the sum form. On their ranges the form is bounded by $hM$ times
the squared Hilbert norm. Algebraic tensors are dense there, and each
factor can in turn be approximated in its one-place form norm by
eventually constant sequences. This proves that such tensors are a
core for the sum form. Symmetrizing preserves convergence and identifies
the slice traces. We may consequently approximate $(u,v)$ by arrays
constant in each place outside a finite set. This argument does not
require traces on intersections of two or more limiting faces.

\paragraph{Return to polynomials.}
The correspondence is already visible with one place and even parity.
If $u_i=v$ outside a finite set $E$, the polynomial
\[
 P(s,y)=v s+\sum_{i\in E}(u_i-v)(y_i^2-l_i)
\]
has $P(S_y,y)=\sum_i u_i(y_i^2-l_i)$ and $\partial_sP=v$.
Its compensated derivatives vanish outside $E$. Thus the constant
tail of a coefficient array is represented by the centered square
sum. For several places, its powers give the corresponding tensor
coefficients, as the following calculation shows.

Choose a finite set $E$ containing $O$ and all exceptional indices of
one of these approximating arrays. A symmetric array constant in each
place outside $E$ is determined by its value $U_j$ for every count
vector $j=(j_i)_{i\in E}$ with $|j|\le h$: the remaining
$r=h-|j|$ places can have arbitrary indices outside $E$, including
coincident ones. Define
\[
 S_{E^c}(s,y)=s-\sum_{i\in E}(y_i^2-l_i),
 \qquad
 M_j(y)=\prod_{i\in E}y_i^{\mathbf1_O(i)+2j_i}.
\]
The degree-$d$ Hermite coefficients of
$M_j(y)S_{E^c}(S_y,y)^r$ have the prescribed counts $j$ in $E$.
For outside counts $k$ with $|k|=r$, their coefficient is
\[
 C_E(j)\,r!\prod_{i\notin E}
              l_i^{k_i}\frac{\sqrt{(2k_i)!}}{k_i!},
 \qquad
 C_E(j)=\prod_{i\in E}
       l_i^{(\mathbf1_O(i)+2j_i)/2}
       \sqrt{(\mathbf1_O(i)+2j_i)!}.
\]
This follows directly from the multinomial expansion, and so includes
repeated outside indices. Multiplying this polynomial by
$U_j c(j)/(r!C_E(j))$, and summing over the finitely many $j$, realizes
exactly the desired top-grade array $u$. Here $c(j)$ is restricted to
$E$. Differentiation in $s$ replaces $S_{E^c}^r$ by
$rS_{E^c}^{r-1}$. Its top coefficients have the same limiting one-place
trace, since $r(r-1)!=r!$. They therefore realize $v$, also by trace
uniqueness.

Every such polynomial pair has a square-summable compensated derivative,
since differentiation after subtracting $2y_i\partial_s P$ leaves only
the finitely many individual-coordinate derivatives. All its moments
are finite; the centered sum $S_y$ has exponential moments in a
neighborhood of zero. Induction on the grade now removes the lower
Hermite grades of these polynomial pairs. More explicitly, their lower
grades satisfy the same summability identities and, by induction,
already belong to the polynomial closure. Subtracting their
approximants shows that the desired top-grade pair lies in that closure.
The fixed-grade approximations and the preceding damping argument prove
the lemma.
\end{proof}

\begin{proof}[Completion of the proof of Proposition~\ref{prop:gradient-domain}]
We can now use Lemma~\ref{lem:gradient-compensated}. Work in the
uniform-angle times product Gaussian norm, equivalent to the localized
conditioned norm. First project onto finitely many spherical harmonics
in $\vartheta$; the projections
converge in the first angular Sobolev norm and commute with the tail
identities. Apply the lemma to each resulting coefficient. This yields
polynomials in $S_y$ and finitely many $y_i$, with smooth angular
coefficients, converging in $f$, $a$, the compensated derivatives, and
the angular derivative.

Multiply these approximants by a smooth cutoff of $S_y$ supported in
a slightly larger radius annulus and equal to one on the original
support. Under multiplication by $\chi(S_y)$ the scalar component
becomes $\chi a+\chi'f$, while the compensated derivative becomes
$\chi G$. This is continuous in all the norms under consideration.
Finally substitute $S_y=C-R^2$. The approximants are now genuine
cylinder functions in $x$, smooth because their first-block radii stay
away from zero. Their Euclidean gradients have exactly the components
in \eqref{eq:source-24}. Ordinary cutoffs in their finitely many other
coordinates make them compactly supported, with convergence in the
same norms. Removing the localizations proves equality of the cylinder
closure and the angular graph.

It remains to identify the kernel. If $G=0$, the restriction of $F$ to
almost every finite initial-block sphere is constant, since the sphere
is connected. Thus $F$ is measurable, modulo null sets, with respect to
the complement of every sufficiently large initial block. This
completed tail field is trivial. To check this directly under $\Pi$,
eliminate one fixed initial block. Its tail law has an integrable
density with respect to product Gaussian measure. Approximating this
density and its product with any bounded observable by cylinder
functions shows that events beyond increasing indices are
asymptotically independent of that observable, uniformly over the
events. An event in every completed tail field is consequently
independent of itself and has probability zero or one. Hence $F$ is
constant. Conversely constants satisfy the angular identities with
$G=0$ and are in the closed domain by what we have proved.

The same characterization applies to $x_a$: it is square integrable
by Proposition~\ref{prop:edge-conditioned-law}, and its angular weak
gradient is $e_a$. This proves the last assertion.
\end{proof}

\subsection{Observation smoothings and their slopes}

We return to the finite SK model and use the quenched subsequential
convention of Remark~\ref{rem:edge-tight-convergence}. Write
$T=n^{2/3}$. Consider functions $f=f_n$ such that
\[
 \mu|f|^2+T\mathcal E(f)\le C,
 \qquad \|f\|_\infty\le Cn^D
\]
for a fixed exponent $D$. The bounds below hold on disorder sets of
arbitrarily high probability, with finite constants depending on the
set but not on $f$. Terms denoted by $o(1)$ absorb the exponentially
small exceptional-path costs multiplied by these polynomial bounds.
The fixed exponent $D$ affects these vanishing errors, not the
constants multiplying the energy.

For the cap observation $h_p$, define
\[
 F_p(h_p)=\mathbb E[f\mid h_p],
 \qquad S_pf(X)=\mathbb E[F_p(h_p)\mid X],
 \qquad
 J_p^f(X)=n^{1/3}\mathbb E[B_p\nabla_{h_p}F_p(h_p)\mid X].
\]
All expectations here are under the ordinary observation law for the
fixed disorder. The vector $J_p^f$ is expressed in physical coordinates.
Recall that $u_{p,a}=(U^t h_p)_a/(p^2n^{1/3})$.

\begin{lemma}[Smoothing estimates]
\label{lem:gradient-smoothing}
For $r\le p\le p_*$ in the range of
Proposition~\ref{prop:variance-main},
\begin{equation}
 Q_\mu\|\nabla_{u_p}F_p\|^2
       \le C T\mathcal E(f)+o(1),
 \label{eq:source-26}
\end{equation}
and
\begin{equation}
 \begin{split}
 \|n^{1/3}(f-S_pf)\|_{L^2(\mu)}
       &\le Cp^{-1}\sqrt{T\mathcal E(f)}+o(1),\\
 \left\|n^{1/3}(D_iS_pf)_i-J_p^f\right\|_{L^2(\mu;\mathbb R^n)}
       &\le Cp\sqrt{T\mathcal E(f)}+o(1).
 \end{split}
 \label{eq:source-27}
\end{equation}
In particular, $J_p^f$ is bounded in $L^2(\mu;\mathbb R^n)$.
\end{lemma}

\begin{proof}
On the active cap,
$\nabla_{h_p}F_p=P_p\Cov(X,f\mid h_p)$. Conditional Cauchy--Schwarz
therefore gives
\[
 |\nabla_{h_p}F_p|^2
 \le\|K_p\|\Var(f\mid h_p).
\]
Combining the covariance bound \eqref{eq:source-6}, the annulus estimate
\eqref{eq:source-7}, and the variance estimate
\eqref{eq:source-16}, and multiplying by the squared change-of-variable
factor $p^4n^{2/3}$, proves \eqref{eq:source-26}. Conditional Jensen and
$\|B_p\|\le Cp^2$ also bound $J_p^f$. The first line of
\eqref{eq:source-27} follows from the same Jensen inequality applied to
$f-F_p(h_p)$.

For the second line, condition on $X$ and write the channel as
$h_p=B_pX+B_p^{1/2}Z$, with $Z$ standard Gaussian on the active space.
Flipping spin $i$ translates $Z$ by $-2X_iB_p^{1/2}e_i$. Expand the
Gaussian translation in Hermite tensors. The degree-one term in the
half-difference is the $i$th component of $J_p^f$. For the higher
degrees, subtract the constant $f(X)$ from the translated function.
The total squared norm of the Hermite coefficients of
$n^{1/3}(F_p(h_p)-f(X))$, averaged over $X$, is at most
\[
 C p^{-2}T\mathcal E(f)+o(1)
\]
by \eqref{eq:source-16}. The degree-$j$ contraction into the vector
indexed by $i$ has operator norm at most
\[
 \frac{2^j}{\sqrt{j!}}\|B_p\|^{1/2}
           \left(\max_i(B_p)_{ii}\right)^{(j-1)/2}.
\]
Indeed its Gram matrix, apart from the displayed scalar and spin signs,
is a Schur power of $B_p$; the Schur product bound gives the estimate.
Since $\max_i(B_p)_{ii}\le\|B_p\|\le Cp^2$, summing degrees $j\ge2$
costs $O(p^2)$ times the coefficient norm. This proves the second line
of \eqref{eq:source-27} in the unweighted vector norm.
\end{proof}

\subsection{Identification of the limiting derivative}

The static convergence of cube cylinder integrals in
Corollary~\ref{cor:observations-static} implies that a bounded $L^2(\mu)$
sequence has a subsequence with a weak cylinder limit $F\in L^2(\Pi)$.
Concretely, this means
$\mu[f\psi(x)]\longrightarrow\Pi[F\psi]$ for a countable dense family
of cylinder tests, and then for all such tests by their $L^2$ bounds.
For fixed small $q>0$, let a bar denote probability averaging in
$d\log p$ over $[q,2q]$.

\begin{proposition}[Limit of the observation slopes]
\label{prop:gradient-observation}
Suppose $f_n$ satisfies the preceding bounds and has weak cylinder
limit $F$. Choose a subsequence on which $T\mathcal E(f_n)$ tends to
its liminf. After further extraction in the quenched convention, there
is a sequence $q_j\downarrow0$ such that the spectral components of
$\overline{J^f}_{q_j}$ have weak cylinder limits $G^{(j)}$ as
$n\to\infty$. Every weak cluster point of $(G^{(j)})$ in
$L^2(\Pi;\ell^2)$ equals $\nabla F$. In particular,
$F$ belongs to the form domain of Proposition~\ref{prop:gradient-domain}
and
\begin{equation}
 \|\nabla F\|_{L^2(\Pi;\ell^2)}
 \le C\left(\liminf_{n\to\infty}T\mathcal E(f_n)\right)^{1/2}.
 \label{eq:source-28}
\end{equation}
The constant is the same for all sequences with the stated size bounds.
\end{proposition}

\begin{proof}
Pass first to a subsequence on which the energy converges to its
liminf. Equation~\eqref{eq:source-26} bounds the weak spectral vector
limits in $L^2(\Pi;\ell^2)$. To identify them, fix $A$ and $\psi$ as
in \eqref{eq:source-23}. We will compare the angular pairing of the
slopes at fixed cap scale to an integration by parts at the smallest
scale $r$.

Define
\[
 T_p^A=Q_\mu[\psi(u_p)\mathcal R_AF_p(h_p)],
\]
where the rotation acts on spectral coordinates. By
\eqref{eq:source-8} and \eqref{eq:source-26},
\[
 \mu[\psi(x)(Ax)\cdot U^tJ_p^f]-T_p^A
       =O\bigl((pn^{1/3})^{-1}\bigr)+o(1).
\]
For the finitely many coordinates involved,
$B_p/p^2=I+O((pn^{1/3})^{-2})$. Replacing $x$ by $u_p$ in the test
then uses Cauchy--Schwarz and \eqref{eq:source-8}; compact support of
the test, or its fourth-moment version, controls the coordinate factors.

\paragraph{Transport between two scales.}
Let $r\le c<p\le2c$, with all indices of $A$ active at scale $c$.
The increment covariance $B_p-B_c$ is scalar on those indices.
Rotation therefore commutes with the Gaussian convolution expressing
the observation bridge. Differentiating its numerator and denominator
gives the exact identity
\begin{equation}
 \mathcal R_AF_c(h_c)
 =\mathbb E\left[
      \mathcal R_AF_p(h_p)
      +(F_p(h_p)-F_c(h_c))\mathcal R_A\phi_p(h_p)
      \,\middle|\,h_c\right].
 \label{eq:source-29}
\end{equation}
In the first term, replacement of $\psi(u_c)$ by $\psi(u_p)$ costs
$O((cn^{1/3})^{-1})+o(1)$, using
\eqref{eq:source-8} and \eqref{eq:source-26}. The other term is bounded
in absolute value by
\begin{equation}
 \|\psi\|_\infty V_c(f)^{1/2}
              \|\mathcal R_A\phi_p\|_{L^2(Q_\mu)}.
 \label{eq:source-30}
\end{equation}
Here $\mathbb E[(F_p-F_c)^2]\le V_c(f)$ follows from the posterior
martingale and conditional variance decomposition.

Starting at each $p\in[q,2q]$, descend by factors of two, using $r$
as the final scale if necessary. The coordinate errors sum to
$O((rn^{1/3})^{-1})=O(n^{-\delta})$. For the score errors,
\eqref{eq:source-16} bounds \eqref{eq:source-30}, up to a fixed tight
factor, by the root mean square of $\mathcal R_A\phi_p$ divided by
$n^{1/3}p$. Average over the starting $p$. In every resulting parent
window, \eqref{eq:source-15} and Cauchy--Schwarz apply, including a
partial final window by enlarging it to its fixed-ratio buffer. Under
the biased law the expected contribution of a window of scale $Q$ is
at most
\[
 C_A\bigl(Q^{\tau'}+k_Q^{-\tau'}\bigr)^{1/2}.
\]
The comparison through $L_\circ\ge1/2$ transfers this estimate to
ordinary disorder integrals. Its sum is bounded by
$C_A(q^{\tau'/2}+k_r^{-\tau'/2})$. Markov's inequality therefore shows
that the averaged $T_p^A$ and $T_r^A$ differ by a quantity tending to
zero in probability as $n\to\infty$ and then $q\downarrow0$. The
estimate is independent of the choice of $f$ within the stipulated
size and energy bounds.

\paragraph{Integration by parts at the smallest scale.}
The observation density in the active $h_r$ coordinates is proportional
to
\[
 \exp\left(\phi_r(h_r)-\frac12h_r^tB_r^{-1}h_r\right).
\]
In $u_r$ coordinates the quadratic coefficient in direction $i$ is
$n^{2/3}r^4/B_{r,i}$. Its pairwise differences are
\begin{equation}
 \delta_{ij}
 =n^{2/3}r^4\frac{d_i-d_j}{B_{r,i}B_{r,j}}
 \longrightarrow g_i-g_j
 \label{eq:source-31}
\end{equation}
for fixed leading indices. The common divergent part of the coefficient
cancels in an angular derivative. Also $\phi_r^{\rm sp}$ is radial on
the cap, so that
$\mathcal R_A\phi_r=\mathcal R_A\log L_r$. Its contribution to the
integration-by-parts identity tends to zero by
\eqref{eq:source-12}, the field-coordinate moments, and the same
biased-to-quenched comparison. Cauchy--Schwarz applies because $F_r$
is bounded in $L^2$ and the relevant score has vanishing $L^2$ norm
after multiplication by the compact cylinder test.

For every function $b(h_r)$,
$Q_\mu[F_r(h_r)b(h_r)]=Q_\mu[f(X)b(h_r)]$. We may therefore use
\eqref{eq:source-8} in the remaining terms of the integration-by-parts
identity and then the weak cylinder convergence of $f$. Equations
\eqref{eq:source-31} and \eqref{eq:source-23} have the same limiting
quadratic coefficient, since adding the scalar $b$ has no angular
effect. We obtain
\[
 T_r^A\longrightarrow
 -\int F\left\{\mathcal R_A\psi
       -\psi\sum_a(Ax)_a(g_a+b)x_a\right\}\,d\Pi.
\]
Together with the preceding scale comparison, this proves that each
weak limit $G$ of the averaged slopes satisfies
\eqref{eq:source-23}. Proposition~\ref{prop:gradient-domain} identifies
it with $\nabla F$, and weak lower semicontinuity in
\eqref{eq:source-26} proves \eqref{eq:source-28}.

Finally, these arguments can be realized simultaneously for the
countably many scales and tests needed in the sequel. Represent the
bounds by tight random constants on the high-probability disorder
sets, and include the error variables in the subsequential realization
of Remark~\ref{rem:edge-tight-convergence}. The errors tend to zero
along a sufficiently rapidly decreasing outer sequence of scales.
On the resulting limit samples, weak Hilbert-space extraction applies
to any of the particular sequences $f_n$ under consideration. This
establishes the stated quenched form of the proposition.
\end{proof}
